\chapter{Fundamentação Teórica}

Este capítulo apresenta os conceitos teóricos que fundamentam o projeto, a implementação e a validação experimental do protótipo de braço robótico EB-15. A revisão abrange desde os princípios gerais de robótica industrial e cinemática de manipuladores até os aspectos específicos de atuadores, sistemas embarcados, integração em rede e métricas de desempenho.

Dessa forma, busca-se fornecer ao leitor a base necessária para compreender as decisões de projeto adotadas e os critérios de avaliação empregados na comparação com o robô industrial de referência. Primeiramente, são abordados os fundamentos de robótica e a classificação de manipuladores industriais, seguidos pela modelagem cinemática por meio de transformações homogêneas e parâmetros de Denavit-Hartenberg.

Em seguida, discutem-se os métodos de planejamento de trajetória e os perfis de velocidade aplicáveis ao controle de movimento. Posteriormente, apresentam-se os atuadores eletromecânicos e os circuitos de acionamento utilizados, bem como as plataformas embarcadas responsáveis pelo controle distribuído. Por fim, são introduzidos os conceitos de Internet das Coisas aplicados à robótica, os mecanismos de segurança lógica e as métricas normalizadas de validação experimental.

\section{Robótica Industrial e Manipuladores}

A robótica industrial tem desempenhado um papel crescente na automação de processos produtivos desde a instalação do primeiro robô manipulador em linhas de montagem na década de 1960. Conforme \citeonline{siciliano2009}, um robô manipulador industrial pode ser definido como um mecanismo programável, composto por uma série de elos rígidos conectados por juntas, cuja função é posicionar e orientar uma ferramenta ou peça no espaço tridimensional. Nesse contexto, a evolução dos microcontroladores e dos componentes de fabricação digital ampliou o acesso a essas tecnologias, permitindo que protótipos funcionais sejam desenvolvidos em ambientes acadêmicos com orçamentos significativamente inferiores aos da indústria \cite{ramirez2022}. Portanto, compreender a estrutura, a classificação e os princípios de funcionamento dos manipuladores constitui o ponto de partida indispensável para o desenvolvimento do protótipo proposto neste trabalho.

\subsection{Robôs Manipuladores -- Classificação e Anatomia}

Os robôs manipuladores industriais são classificados de acordo com a configuração cinemática de suas juntas e elos, que determina o espaço de trabalho e a destreza do mecanismo. Segundo \citeonline{craig2018}, as configurações mais comuns incluem os robôs articulados (antropomórficos), SCARA (\textit{Selective Compliance Assembly Robot Arm}), cartesianos (retangulares) e Delta (paralelos), cada qual adequado a requisitos específicos de alcance, velocidade e carga. De modo complementar, \citeonline{spong2020} destacam que a quantidade de graus de liberdade (DOF — \textit{Degrees of Freedom}) de um manipulador define sua capacidade de posicionamento e orientação no espaço, sendo necessários ao menos seis graus de liberdade para que o efetuador final atinja uma pose arbitrária em três dimensões. Sendo assim, a cadeia cinemática de um robô manipulador — entendida como a sequência de juntas rotativas e prismáticas que conectam a base ao efetuador — constitui o parâmetro estrutural que governa tanto o volume de trabalho quanto a complexidade do controle.

Nessa perspectiva, o protótipo EB-15 objeto deste trabalho enquadra-se na categoria de robô articulado com seis juntas rotativas, configuração semelhante à adotada por manipuladores industriais consolidados como o UR-10 da Universal Robots. Conforme \citeonline{siciliano2009}, robôs articulados de seis eixos são amplamente utilizados em aplicações de soldagem, montagem e manipulação de materiais, precisamente por oferecerem a combinação de alcance espacial e versatilidade de orientação requerida por essas tarefas. Todavia, enquanto plataformas industriais empregam servomotores de alta precisão com realimentação de posição, o protótipo de baixo custo fundamenta-se em motores de passo sem encoders, o que impõe limitações experimentais discutidas nas seções subsequentes. Por conseguinte, posicionar ambos os robôs dentro dessa taxonomia permite estabelecer o quadro comparativo que orienta toda a metodologia experimental do trabalho.

Cabe destacar que o UR-10 (Universal Robots) é um robô colaborativo de seis graus de liberdade com capacidade de carga de 10 kg, alcance de 1300 mm, repetibilidade de $\pm$0,1 mm e controle em malha fechada com encoders absolutos em todas as juntas \cite{universalrobots2023}. Dessa forma, suas especificações representam o patamar industrial de referência contra o qual o protótipo EB-15 será avaliado experimentalmente.

\subsection{Projeto Mecânico de Braços Robóticos de Baixo Custo}


O avanço da impressão 3D e a consolidação de comunidades de hardware aberto viabilizaram o surgimento de projetos de braços robóticos acessíveis, destinados a fins educacionais e de pesquisa. De acordo com \citeonline{gonzalezgomez2012}, plataformas robóticas \textit{open-source} baseadas em peças impressas e componentes padronizados reduzem drasticamente o custo de entrada para laboratórios e universidades, sem comprometer a funcionalidade necessária para experimentação em cinemática e controle. Desse modo, projetos como o Thor, o BCN3D Moveo e o próprio EB-15 demonstram que é possível construir manipuladores de seis graus de liberdade com investimentos da ordem de centenas de reais, em contraste com os milhares de dólares exigidos por equipamentos industriais.

Não obstante, a adoção de materiais poliméricos impressos e redutores compactos introduz compromissos significativos em termos de rigidez estrutural, tolerância dimensional e repetibilidade mecânica. Segundo \citeonline{niku2020}, a deflexão dos elos sob carga e a folga (\textit{backlash}) nos elementos de transmissão representam as principais fontes de erro posicional em manipuladores leves de baixo custo. Por conseguinte, a viabilidade técnica de um protótipo dessa natureza depende não apenas da adequação do projeto mecânico aos requisitos operacionais, mas também da capacidade do sistema de controle em compensar ou tolerar essas imperfeições dentro de limites aceitáveis para a aplicação pretendida.

\subsection{Reduções Mecânicas -- Engrenagens Planetárias}

As caixas de redução planetárias constituem um dos elementos mecânicos fundamentais na transmissão de torque em braços robóticos, sendo amplamente adotadas por oferecerem alta relação de redução em um volume compacto. Conforme \citeonline{shigley2011}, um trem epicicloidal é composto por uma engrenagem solar central, um conjunto de engrenagens planetárias (satélites) e um anel externo (coroa), cuja configuração permite distribuir a carga entre múltiplos pontos de contato simultâneos. Dessa forma, a capacidade de torque do conjunto é significativamente superior à de uma redução por engrenagens comuns de mesmo tamanho, característica particularmente relevante para juntas robóticas onde o espaço disponível é limitado.

Ademais, a relação de redução alcançada em um estágio planetário pode variar tipicamente entre 3:1 e 10:1, podendo ser ampliada em configurações de múltiplos estágios. Segundo \citeonline{siciliano2009}, a eficiência mecânica de redutores planetários situa-se geralmente entre 90\% e 97\%, enquanto a folga angular (\textit{backlash}) — embora superior à de redutores harmônicos — permanece em níveis aceitáveis para aplicações educacionais e de prototipagem. Portanto, a utilização de caixas planetárias nas juntas do EB-15 representa um compromisso adequado entre custo, compacidade e desempenho de torque, ao passo que a folga mecânica residual constitui um fator limitante a ser considerado na análise experimental de precisão posicional.

\section{Cinemática de Manipuladores}

A cinemática de manipuladores constitui o ramo da mecânica que estuda o movimento de cadeias articuladas sem considerar as forças que o produzem. Conforme \citeonline{craig2018}, a modelagem cinemática de um robô manipulador envolve duas formulações complementares: a cinemática direta, que determina a pose do efetuador a partir dos ângulos das juntas, e a cinemática inversa, que calcula os ângulos necessários para atingir uma pose desejada. Ambas as formulações dependem de uma representação matemática rigorosa das relações espaciais entre os elos da cadeia, o que constitui o objeto das subseções a seguir.

\subsection{Representação Espacial -- Matrizes de Transformação Homogênea}

A descrição da posição e orientação de corpos rígidos no espaço tridimensional fundamenta-se no uso de matrizes de transformação homogênea, que combinam rotação e translação em uma única representação matricial 4$\times$4. De acordo com \citeonline{siciliano2009}, uma matriz de transformação homogênea $T \in \mathbb{R}^{4 \times 4}$ é composta por uma submatriz de rotação $R \in \mathbb{R}^{3 \times 3}$, um vetor de translação $\mathbf{p} \in \mathbb{R}^{3}$ e uma linha de perspectiva $[0 \ 0 \ 0 \ 1]$, conforme a estrutura apresentada na Equação \ref{eq:transf_homogenea}. Dessa forma, a composição de transformações entre sistemas de coordenadas consecutivos é realizada por meio da multiplicação matricial, permitindo mapear a posição de qualquer ponto da cadeia cinemática em relação ao referencial da base.

\begin{equation}
T = \begin{bmatrix} R_{3\times3} & \mathbf{p}_{3\times1} \\ 0_{1\times3} & 1 \end{bmatrix}
\label{eq:transf_homogenea}
\end{equation}

Sob o mesmo ponto de vista, \citeonline{craig2018} ressalta que a principal vantagem dessa representação reside na possibilidade de encadear múltiplas transformações de forma sequencial: a pose do efetuador final em relação à base é obtida pela multiplicação das matrizes de transformação de cada elo, isto é, $T_0^n = T_0^1 \cdot T_1^2 \cdot \ldots \cdot T_{n-1}^n$. Sendo assim, as matrizes de transformação homogênea constituem a ferramenta matemática central tanto para o cálculo da cinemática direta quanto para a verificação de posições no espaço de trabalho do manipulador.

\subsection{Parâmetros de Denavit-Hartenberg (DH)}

A convenção de Denavit-Hartenberg fornece um método sistemático para atribuir sistemas de coordenadas a cada elo de um manipulador serial e descrever a relação geométrica entre juntas consecutivas por meio de apenas quatro parâmetros. Conforme \citeonline{craig2018}, os parâmetros DH de cada junta são: o comprimento do elo $a_i$ (distância entre os eixos das juntas $i$ e $i+1$ ao longo do eixo $x_i$), a torção do elo $\alpha_i$ (ângulo entre os eixos das juntas $i$ e $i+1$ em torno do eixo $x_i$), o deslocamento da junta $d_i$ (distância ao longo do eixo $z_{i-1}$) e o ângulo da junta $\theta_i$ (rotação em torno do eixo $z_{i-1}$). Dessa forma, a matriz de transformação entre dois elos consecutivos pode ser expressa como uma função desses quatro parâmetros, conforme a Equação \ref{eq:dh_matrix}.

\begin{equation}
A_i = \begin{bmatrix}
\cos\theta_i & -\sin\theta_i \cos\alpha_i & \sin\theta_i \sin\alpha_i & a_i \cos\theta_i \\
\sin\theta_i & \cos\theta_i \cos\alpha_i & -\cos\theta_i \sin\alpha_i & a_i \sin\theta_i \\
0 & \sin\alpha_i & \cos\alpha_i & d_i \\
0 & 0 & 0 & 1
\end{bmatrix}
\label{eq:dh_matrix}
\end{equation}

Além disso, \citeonline{spong2020} enfatizam que a principal virtude da convenção DH é a redução do número de parâmetros necessários para descrever completamente uma cadeia cinemática serial: enquanto a descrição genérica de uma transformação rígida requer seis parâmetros independentes (três para translação e três para rotação), a convenção DH explora as restrições geométricas inerentes às juntas para reduzir esse número a quatro. Portanto, a tabela DH de um manipulador constitui uma representação compacta e padronizada de sua geometria, sendo utilizada como entrada para os algoritmos de cinemática direta e inversa implementados no firmware do protótipo EB-15.

\subsection{Cinemática Direta}

A cinemática direta (\textit{Forward Kinematics} — FK) consiste no cálculo da posição e orientação do efetuador final de um manipulador a partir dos valores conhecidos das variáveis de junta. De acordo com \citeonline{siciliano2009}, o problema da cinemática direta é resolvido de forma determinística pela multiplicação sequencial das matrizes de transformação DH de cada junta, resultando na matriz de transformação total $T_0^n$ que relaciona o referencial do efetuador ao referencial da base. Desse modo, para um manipulador de seis graus de liberdade, a pose do \textit{Tool Center Point} (TCP) é obtida conforme a Equação \ref{eq:fk}.

\begin{equation}
T_0^6 = A_1(\theta_1) \cdot A_2(\theta_2) \cdot A_3(\theta_3) \cdot A_4(\theta_4) \cdot A_5(\theta_5) \cdot A_6(\theta_6)
\label{eq:fk}
\end{equation}

Sob o mesmo ponto de vista, \citeonline{corke2017} destaca que a extração de grandezas utilizáveis a partir da matriz resultante — como as coordenadas cartesianas $(x, y, z)$ do efetuador e os ângulos de orientação (comumente representados por Euler na forma \textit{roll-pitch-yaw}) — é uma etapa computacionalmente direta, embora requeira atenção a singularidades na conversão de matrizes de rotação para representações angulares. Por conseguinte, a cinemática direta constitui uma operação fundamental tanto para a telemetria em tempo real da pose do robô quanto para a verificação de colisão no planejador de trajetórias.

\subsection{Cinemática Inversa}

A cinemática inversa (\textit{Inverse Kinematics} -- IK) aborda o problema complementar à cinemática direta: dada uma pose cartesiana desejada para o efetuador final, determinar o conjunto de variáveis de junta que a produz. Segundo \citeonline{craig2018}, este problema é intrinsecamente mais complexo que a cinemática direta por três razões fundamentais: a existência de múltiplas soluções válidas para uma mesma pose, a possibilidade de que a pose desejada seja inatingível e a ocorrência de singularidades cinemáticas que tornam certas configurações matematicamente indeterminadas. Dessa forma, a escolha do método de resolução depende tanto da geometria do manipulador quanto dos recursos computacionais disponíveis.

Nesse contexto, \citeonline{siciliano2009} classificam os métodos de solução da cinemática inversa em duas categorias principais: métodos analíticos (ou de forma fechada), que exploram a geometria específica do manipulador para derivar equações algébricas diretas, e métodos numéricos iterativos, que convergem para uma solução por meio de algoritmos como o Jacobiano transposto ou o método de Newton-Raphson. Sendo assim, os métodos analíticos oferecem maior eficiência computacional quando aplicáveis, mas dependem de geometrias específicas que nem sempre são encontradas em manipuladores genéricos.

De modo complementar, \citeonline{lynch2017} destacam que métodos numéricos como a descida cíclica de coordenadas (CCD -- \textit{Cyclic Coordinate Descent}) são particularmente adequados para implementação em plataformas de recursos limitados, pois podem ser interrompidos após um número fixo de iterações, oferecendo soluções aproximadas com erro controlado. Por conseguinte, a implementação da cinemática inversa no protótipo EB-15 emprega uma abordagem iterativa, cujas características e limitações são detalhadas no capítulo de metodologia.

\subsection{Tool Center Point (TCP)}

O conceito de \textit{Tool Center Point} (TCP) designa o ponto de referência localizado na extremidade do efetuador final do robô, em relação ao qual são definidos os comandos de posicionamento cartesiano. Conforme \citeonline{niku2020}, o TCP representa a origem do sistema de coordenadas da ferramenta e pode ser configurado pelo operador de acordo com a geometria do dispositivo acoplado ao último elo, seja uma garra, um sensor ou uma ferramenta de processo. Sendo assim, a distinção entre o referencial da base, o referencial do efetuador e o referencial da ferramenta (TCP) é fundamental para a correta interpretação dos comandos de movimentação e para a exibição de dados de telemetria nas interfaces de operação do sistema.

\section{Planejamento de Trajetória e Controle de Movimento}

O planejamento de trajetória define a lei temporal segundo a qual um robô manipulador percorre o caminho entre uma configuração inicial e uma configuração final. Segundo \citeonline{siciliano2009}, a distinção entre caminho (sequência geométrica de pontos) e trajetória (caminho associado a uma lei temporal) é fundamental, pois um mesmo caminho pode ser percorrido com perfis de velocidade distintos, resultando em comportamentos dinâmicos e esforços mecânicos significativamente diferentes. Ademais, \citeonline{lynch2017} ressaltam que a escolha do espaço de planejamento — espaço de juntas ou espaço cartesiano — influencia tanto a previsibilidade da trajetória do efetuador quanto a carga computacional exigida do controlador. Sendo assim, esta seção apresenta os conceitos de interpolação e geração de perfis de velocidade aplicados ao controle de movimento do protótipo EB-15.

\subsection{Planejamento no Espaço de Juntas vs.\ Espaço Cartesiano}

O planejamento de trajetória pode ser realizado no espaço de juntas (\textit{joint-space}) ou no espaço cartesiano (\textit{task-space}), cada abordagem apresentando características distintas quanto à previsibilidade do movimento e ao custo computacional. Conforme \citeonline{craig2018}, no planejamento em espaço de juntas (designado MoveJ em robôs industriais), as variáveis de junta são interpoladas diretamente entre os valores inicial e final, resultando em um movimento suave e eficiente do ponto de vista articular, porém cuja trajetória do efetuador no espaço tridimensional não segue necessariamente uma linha reta. Em contrapartida, o planejamento em espaço cartesiano (designado MoveL) interpola a posição e a orientação do TCP ao longo de uma linha reta no espaço operacional, exigindo a resolução da cinemática inversa em cada ponto intermediário da trajetória.

Sob o mesmo ponto de vista, \citeonline{lynch2017} destacam que a escolha entre os dois modos de planejamento depende dos requisitos da tarefa: operações de posicionamento ponto a ponto favorecem o MoveJ por sua simplicidade e previsibilidade articular, enquanto tarefas que exigem trajetórias lineares do efetuador — como soldagem ou transporte de objetos — requerem o MoveL apesar de seu maior custo computacional. Por conseguinte, o protótipo EB-15 implementa ambas as modalidades, utilizando MoveJ para movimentação geral e MoveL para deslocamentos lineares que exigem controle da trajetória do TCP.

\subsection{Perfis de Velocidade -- Trapezoidal e S-Curve}

A aplicação de perfis de velocidade controlados é essencial para garantir movimentos suaves e mecanicamente seguros em robôs manipuladores, evitando variações bruscas de aceleração que podem provocar vibrações excessivas e desgaste prematuro dos componentes. De acordo com \citeonline{siciliano2009}, o perfil trapezoidal de velocidade constitui a abordagem mais amplamente utilizada em controle de movimento, sendo composto por três fases: aceleração constante até a velocidade de cruzeiro, manutenção da velocidade constante e desaceleração constante até a parada. Dessa forma, o perfil trapezoidal oferece um compromisso favorável entre simplicidade de implementação e suavidade de movimento, sendo particularmente adequado para sistemas embarcados com capacidade computacional limitada.

Adicionalmente, \citeonline{lynch2017} apresentam o perfil S-curve como uma evolução do trapezoidal, na qual as transições entre as fases de aceleração e velocidade constante são suavizadas por meio de segmentos de aceleração crescente e decrescente (\textit{jerk} controlado), resultando em esforços mecânicos significativamente menores nos atuadores e nas estruturas. Todavia, a implementação de perfis S-curve exige maior capacidade de processamento e maior número de parâmetros de configuração, fatores que podem ser restritivos em plataformas embarcadas de baixo custo. Portanto, o perfil trapezoidal foi adotado como estratégia primária de planejamento no protótipo EB-15, embora a arquitetura do firmware permita a futura incorporação de perfis mais sofisticados.

\subsection{Segmentação Temporal de Trajetórias}

A segmentação temporal consiste na divisão de uma trajetória completa em micro-segmentos de duração fixa, técnica empregada para viabilizar a execução coordenada de movimentos em arquiteturas de controle distribuído. Conforme \citeonline{siciliano2009}, a interpolação em tempo real de trajetórias de múltiplos eixos requer que o controlador gere posições intermediárias a uma taxa fixa — denominada taxa de interpolação — de modo que cada segmento temporal contenha os deslocamentos incrementais para todos os eixos do manipulador. Dessa forma, a execução de um movimento longo é substituída pela execução sequencial de segmentos curtos, cada qual com duração pré-determinada e deslocamentos proporcionais à velocidade desejada.

Além disso, \citeonline{corke2017} destaca que a segmentação temporal é particularmente relevante em arquiteturas nas quais o planejamento e a execução são realizados por processadores distintos, pois os segmentos podem ser transmitidos por meio de filas (\textit{buffers} circulares) que desacoplam a produção e o consumo de dados de movimento. Portanto, essa técnica é empregada no protótipo EB-15 para coordenar os eixos J1–J3 (controlados localmente pelo ESP32-S3) e os eixos J4–J6 (controlados pelo Arduino Uno), garantindo sincronização temporal entre os dois controladores por meio de segmentos de duração uniforme.

\subsection{Coordenação Multi-Eixo}

A execução simultânea de movimentos em múltiplos eixos requer algoritmos que coordenem a geração de pulsos de forma proporcional ao deslocamento de cada eixo, garantindo que todos os eixos concluam seus respectivos percursos no mesmo intervalo de tempo. Conforme \citeonline{craig2018}, algoritmos de interpolação derivados do método de Bresenham -- originalmente desenvolvido para rasterização de linhas em computação gráfica -- são amplamente empregados em controladores de movimento para escalonar a geração de pulsos entre eixos com deslocamentos distintos, assegurando a suavidade e a simultaneidade do movimento resultante. No protótipo EB-15, essa técnica é empregada pelo Arduino Uno para coordenar os eixos J4--J6, escalonando a geração de passos proporcionalmente ao deslocamento máximo entre os três eixos.

\section{Atuadores e Eletrônica de Potência}

Os atuadores constituem os elementos responsáveis pela conversão de sinais elétricos em movimento mecânico nas juntas de um robô manipulador. Conforme \citeonline{hughes2019}, a seleção do tipo de motor e do circuito de acionamento é determinante tanto para o desempenho cinemático quanto para a precisão posicional do sistema, uma vez que as características dinâmicas do atuador — torque, velocidade, resolução e modo de controle — definem diretamente as capacidades e limitações do manipulador. Sendo assim, esta seção apresenta os fundamentos dos motores de passo e dos drivers de potência utilizados no protótipo EB-15.

\subsection{Motores de Passo -- Princípios e Características}

Os motores de passo são atuadores eletromecânicos que convertem pulsos elétricos discretos em incrementos angulares precisos, característica que os torna particularmente adequados para aplicações de posicionamento sem necessidade de sensor de realimentação. Segundo \citeonline{acarnley2002}, o princípio de funcionamento baseia-se na comutação sequencial de bobinas no estator, que induz o alinhamento do rotor em posições angulares discretas denominadas passos, cuja resolução típica varia entre 0,9° e 1,8° por passo completo. Desse modo, a posição angular do motor pode ser controlada de forma incremental pela contagem de pulsos enviados ao driver, dispensando sensores de posição em aplicações de malha aberta.

Adicionalmente, a operação em modos de \textit{micro-stepping} permite subdividir cada passo completo em frações de até 1/32, aumentando a resolução angular efetiva e reduzindo a vibração mecânica do rotor. Conforme \citeonline{hughes2019}, motores da categoria NEMA 17 — especificação que designa o tamanho da face de montagem de 42,3 mm — oferecem tipicamente torques de retenção (\textit{holding torque}) entre 0,2 Nm e 0,5 Nm, valores adequados para as juntas de manipuladores de pequeno porte quando combinados com redutores de engrenagens. Todavia, a curva de torque versus velocidade desses motores apresenta uma queda acentuada em rotações elevadas, impondo uma limitação direta à velocidade máxima de movimentação das juntas do protótipo.

\subsection{Drivers de Motor de Passo -- TB6600}

Os drivers de motor de passo são circuitos de potência responsáveis por converter sinais lógicos de baixo nível em correntes de acionamento nas bobinas do motor, implementando o controle de corrente constante necessário para operação em micro-stepping. De acordo com \citeonline{acarnley2002}, drivers baseados em ponte H completa com modulação de corrente por \textit{chopper} regulam a corrente nas bobinas independentemente da velocidade e da indutância do motor, garantindo a manutenção do torque em uma faixa mais ampla de velocidades. Dessa forma, a interface de controle desses drivers é reduzida a dois sinais digitais — pulso (PUL) e direção (DIR) —, o que simplifica significativamente a integração com microcontroladores.

Sob o mesmo ponto de vista, o driver TB6600 adota essa arquitetura e permite a configuração de micro-stepping (1, 2, 4, 8 e 16 micro-passos por passo completo) e da corrente máxima por meio de chaves DIP físicas no módulo. Conforme \citeonline{hughes2019}, a resolução de micro-stepping selecionada impacta diretamente a relação entre o número de pulsos e o deslocamento angular resultante (\textit{steps/degree}), parâmetro essencial para a calibração do sistema de controle. Portanto, a configuração adequada dos drivers TB6600 constitui um dos primeiros procedimentos de ajuste do protótipo EB-15, influenciando tanto a suavidade de movimento quanto a precisão posicional alcançável.

\subsection{Controle em Malha Aberta vs. Malha Fechada}

Sistemas de controle de posição podem operar em malha aberta (\textit{open-loop}), nos quais a posição é inferida pela contagem de pulsos enviados ao motor, ou em malha fechada (\textit{closed-loop}), nos quais sensores de realimentação — como encoders ópticos ou magnéticos — fornecem medição direta da posição real do eixo. Segundo \citeonline{ogata2010}, o controle em malha aberta pressupõe uma relação determinística entre comando e resposta, sendo vulnerável a perturbações externas que rompam essa relação, como perda de passos em motores de passo submetidos a cargas excessivas ou acelerações superiores à capacidade do motor. Em contrapartida, o controle em malha fechada detecta e corrige desvios entre a posição comandada e a posição real, resultando em maior robustez e precisão.

Não obstante, a implementação de malha fechada requer a adição de sensores de posição em cada junta, a complexificação do firmware para processamento de sinais de realimentação em tempo real e o projeto de controladores compensadores (como PID), resultando em um aumento significativo de custo e complexidade do sistema. Conforme \citeonline{niku2020}, a escolha entre as duas abordagens constitui um dos compromissos fundamentais no projeto de robôs de baixo custo: o controle em malha aberta reduz custo e complexidade em troca de menor robustez à perda de passos e menor precisão absoluta. Sendo assim, o protótipo EB-15 adota controle em malha aberta como estratégia primária, e a quantificação das limitações decorrentes dessa escolha constitui um dos objetivos específicos da validação experimental.

\section{Sistemas Embarcados para Controle Robótico}

Os sistemas embarcados fornecem a plataforma computacional sobre a qual o firmware de controle do robô é executado, integrando capacidades de processamento, comunicação e gerenciamento de periféricos em um único dispositivo. Conforme \citeonline{white2024}, o projeto de sistemas embarcados para controle de movimento em tempo real impõe requisitos específicos de temporização, gerenciamento de interrupções e proteção de dados compartilhados que diferem substancialmente da programação de aplicações convencionais. Dessa forma, esta seção apresenta as características dos microcontroladores utilizados no protótipo EB-15 e os fundamentos de programação embarcada que sustentam a implementação do firmware de controle.

\subsection{Microcontroladores -- ESP32-S3 e ATmega328P}

O protótipo EB-15 emprega dois microcontroladores com capacidades complementares: o ESP32-S3, da Espressif Systems, como controlador supervisor, e o ATmega328P (Arduino Uno), da Microchip Technology, como executor subordinado. Segundo \citeonline{espressif2023}, o ESP32-S3 é um \textit{System-on-Chip} (SoC) baseado em dois núcleos Xtensa LX7 operando a até 240 MHz, equipado com 512 KB de SRAM, conectividade Wi-Fi 802.11 b/g/n e Bluetooth 5 (LE), além de ricos periféricos de temporização e comunicação. Dessa forma, suas capacidades de processamento e conectividade o tornam adequado para funções de alto nível como planejamento de trajetória, cálculo de cinemática direta e inversa, hospedagem de servidor HTTP e gerenciamento de API REST.

Em contrapartida, \citeonline{atmel2018} descreve o ATmega328P como um microcontrolador de 8 bits da família AVR com 32 KB de memória Flash, 2 KB de SRAM e frequência de operação de 16 MHz, recursos significativamente mais limitados. Todavia, a simplicidade arquitetural do ATmega328P e o determinismo temporal de seu \textit{loop} de execução o tornam adequado para a tarefa específica de execução de segmentos de movimento: gerar pulsos de controle para os eixos J4–J6 com temporização precisa e gerenciar uma fila circular de segmentos recebidos do ESP32-S3. Portanto, a divisão de responsabilidades entre os dois controladores explora os pontos fortes de cada plataforma, configurando uma arquitetura mestre-escravo distribuída.

\subsection{Interrupções e Temporização em Sistemas de Tempo Real}

A geração de pulsos de controle para motores de passo em frequências da ordem de quilohertz exige mecanismos de temporização por hardware que operem de forma independente do \textit{loop} principal do firmware. De acordo com \citeonline{white2024}, rotinas de serviço de interrupção (ISR — \textit{Interrupt Service Routines}) vinculadas a \textit{timers} de hardware permitem a execução de código em intervalos precisos e determinísticos, garantindo que a geração de pulsos não seja afetada pela latência variável das funções executadas no \textit{loop} principal. Desse modo, o ESP32-S3 utiliza um \textit{timer} de hardware configurado para disparar interrupções a cada 50 microssegundos (frequência de 20 kHz), dentro das quais os pulsos para os motores J1–J3 são gerados com temporização precisa.

Adicionalmente, \citeonline{espressif2023} destaca que o compartilhamento de dados entre o contexto da ISR e o \textit{loop} principal requer o uso de seções críticas (\textit{critical sections}) ou mecanismos de exclusão mútua (\textit{mutexes}), a fim de evitar condições de corrida que poderiam corromper variáveis de estado do motor. Portanto, a correta implementação do gerenciamento de interrupções e da proteção de dados compartilhados é um requisito fundamental para a estabilidade e a confiabilidade do controle de movimento em tempo real.

\subsection{Comunicação UART e Protocolo Binário}

A comunicação entre o ESP32-S3 e o Arduino Uno é realizada por meio de uma interface UART (\textit{Universal Asynchronous Receiver-Transmitter}), sobre a qual foi implementado um protocolo binário estruturado com mecanismos de integridade e confirmação. Conforme \citeonline{axelson2007}, protocolos seriais para sistemas embarcados tipicamente empregam bytes delimitadores de início e fim (\textit{framing}), campos de comprimento, verificação de integridade por CRC (\textit{Cyclic Redundancy Check}) e mensagens de confirmação (ACK) e rejeição (NACK), formando um enlace de comunicação robusto mesmo sobre uma camada física simples como a UART.

\subsection{Arquitetura Mestre-Escravo}

A arquitetura de controle distribuído do protótipo EB-15 segue o paradigma mestre-escravo, no qual o ESP32-S3 atua como nó decisor — responsável pelo planejamento de trajetórias, pela resolução de cinemática e pela comunicação com o operador — enquanto o Arduino Uno atua como nó executor, limitando-se a receber e executar segmentos de movimento. Conforme \citeonline{white2024}, essa separação de responsabilidades permite que cada processador opere dentro de suas capacidades, ao passo que mecanismos de \textit{watchdog} e \textit{heartbeat} periódico garantem a detecção de falhas de comunicação.

\subsection{Persistência de Dados em Microcontroladores}

O armazenamento de parâmetros de configuração que devem sobreviver a ciclos de energia é realizado por meio de técnicas de persistência não volátil integradas aos microcontroladores. Conforme \citeonline{espressif2023}, o ESP32-S3 oferece dois mecanismos principais: o NVS (\textit{Non-Volatile Storage}), que armazena pares chave-valor em partições da memória Flash interna, e o SPIFFS (\textit{SPI Flash File System}), que implementa um sistema de arquivos leve para armazenamento de dados estruturados como arquivos JSON.

\section{IoT e Interface de Operação}

A integração de dispositivos físicos a redes de comunicação constitui um dos pilares da evolução tecnológica denominada Internet das Coisas (IoT), cuja aplicação em sistemas robóticos amplia as possibilidades de monitoramento, controle remoto e interoperabilidade. Conforme \citeonline{atzori2010}, o paradigma IoT pressupõe a capacidade de objetos físicos interagirem com a infraestrutura de rede de forma autônoma, disponibilizando dados de estado e aceitando comandos por meio de protocolos padronizados. Dessa forma, esta seção apresenta os conceitos de IoT e de interfaces de operação que fundamentam a conectividade e a interface web do protótipo EB-15.

\subsection{Internet das Coisas — Conceitos e Aplicações na Robótica}

A Internet das Coisas (IoT) pode ser definida como o paradigma tecnológico no qual objetos do mundo físico são dotados de capacidade de comunicação em rede, permitindo a coleta de dados, o acionamento remoto e a integração com sistemas de informação. Segundo \citeonline{atzori2010}, a convergência entre sensores, microcontroladores com conectividade sem fio e protocolos de comunicação padronizados viabilizou a expansão da IoT para além de aplicações domésticas e de consumo, alcançando domínios como manufatura, logística e robótica experimental. De modo complementar, \citeonline{gubbi2013} situam a IoT como elemento central da quarta revolução industrial (Indústria 4.0), na qual a interconexão de máquinas, sensores e sistemas de informação permite a criação de ambientes de produção inteligentes e adaptativos.

Nesse cenário, a aplicação de conceitos IoT em sistemas robóticos permite que um braço manipulador seja monitorado e comandado remotamente por meio de dispositivos conectados à mesma rede local, sem a necessidade de cabos dedicados ou consoles proprietários. Conforme \citeonline{schwab2016}, a Indústria 4.0 se caracteriza pela fusão de tecnologias digitais, físicas e biológicas, e a conectividade de robôs é um dos exemplos mais tangíveis dessa integração. Portanto, a integração IoT do protótipo EB-15 — que hospeda uma API REST e uma interface web diretamente no controlador ESP32-S3 — insere-se nesse contexto de democratização do acesso a funcionalidades de controle e supervisão anteriormente restritas a sistemas industriais proprietários.

\subsection{API REST e Interface Web em Sistemas Embarcados}

O estilo arquitetural REST (\textit{Representational State Transfer}), proposto originalmente por \citeonline{fielding2000}, define princípios para o projeto de interfaces de comunicação em sistemas distribuídos baseados em HTTP, nos quais os recursos do sistema são representados como URIs e manipulados por meio dos verbos padrão do protocolo (GET, POST, PUT, DELETE). Quando aplicados a dispositivos embarcados com recursos limitados, esses princípios permitem a criação de interfaces programáticas leves que dispensam \textit{frameworks} complexos. De modo análogo, interfaces web embarcadas — nas quais o microcontrolador serve páginas HTML/CSS/JavaScript diretamente a partir de sua memória interna — representam uma abordagem de baixo custo para a implementação de painéis supervisórios (IHM) que substituem consoles proprietários em protótipos experimentais.

\section{Segurança Lógica em Sistemas Robóticos}

A operação segura de robôs manipuladores requer a implementação de mecanismos de proteção que previnam danos ao equipamento, ao ambiente e aos operadores. Conforme \citeonline{siciliano2009}, as camadas de segurança em sistemas robóticos incluem tanto proteções físicas (barreiras, chaves de fim de curso, limitadores de torque) quanto proteções lógicas implementadas por software, sendo estas últimas especialmente relevantes em protótipos experimentais que podem carecer da robustez mecânica de plataformas industriais.

\subsection{Segurança Lógica -- Soft Limits, Watchdog e Parada de Emergência}

Os limites de software (\textit{soft limits}) constituem restrições programáticas que impedem que as variáveis de junta excedam faixas angulares pré-definidas, protegendo a estrutura mecânica contra movimentos que ultrapassem os limites físicos de cada articulação. Segundo \citeonline{niku2020}, a verificação de limites de junta deve ser realizada antes do início de cada movimento planejado e, idealmente, em cada ponto intermediário da trajetória, de modo a detectar violações que possam surgir durante a interpolação. Desse modo, a combinação de limites de software com limites de velocidade máxima por eixo cria uma primeira camada de proteção contra movimentos potencialmente destrutivos.

Além disso, o mecanismo de \textit{watchdog} de comunicação monitora continuamente a integridade do enlace serial entre os controladores, acionando uma parada automática caso a comunicação seja interrompida por um período superior a um intervalo de segurança pré-definido. Conforme \citeonline{white2024}, \textit{watchdog timers} são amplamente utilizados em sistemas embarcados críticos para detectar travamentos de software e falhas de comunicação, garantindo que o sistema transite para um estado seguro mesmo na ausência de um operador. Portanto, o protótipo EB-15 implementa um conjunto integrado de mecanismos de segurança lógica — incluindo \textit{soft limits}, limites de velocidade, \textit{watchdog} serial, transições de estado protegidas e validação de \textit{payloads} na API — que representam uma abordagem defensiva ao controle de um manipulador experimental sem as redundâncias de segurança típicas de robôs industriais.

\section{Validação Experimental em Robótica}

A avaliação quantitativa do desempenho de robôs manipuladores requer a adoção de métricas padronizadas que permitam comparações objetivas entre plataformas distintas. A norma ISO 9283 \cite{iso9283} estabelece os critérios de desempenho e os métodos de teste internacionalmente reconhecidos para a caracterização de robôs industriais, fornecendo definições precisas e procedimentos reprodutíveis para cada métrica. Conforme \citeonline{siciliano2009}, essas métricas constituem a base para certificação e comparação de plataformas robóticas em contextos acadêmicos e industriais.

\subsection{Métricas de Desempenho -- Precisão, Repetibilidade e Tempo de Resposta}

A norma ISO 9283 \cite{iso9283} define as métricas fundamentais de desempenho para robôs manipuladores industriais, entre as quais se destacam a precisão posicional (\textit{pose accuracy}), a repetibilidade posicional (\textit{pose repeatability}) e o tempo de estabilização (\textit{position stabilization time}). Segundo \citeonline{craig2018}, a precisão posicional refere-se ao desvio entre a posição comandada e a posição efetivamente alcançada pelo efetuador, sendo influenciada por erros de calibração, folgas mecânicas e limitações do sistema de controle. Em contrapartida, a repetibilidade mede a dispersão entre múltiplas execuções de um mesmo comando, indicando a consistência do sistema independentemente de seu erro absoluto.

De modo complementar, \citeonline{spong2020} destacam que a repetibilidade tende a ser significativamente melhor que a precisão absoluta em robôs manipuladores, uma vez que erros sistemáticos (como imprecisões na tabela DH) afetam a precisão de forma consistente, mas não comprometem a consistência entre repetições. Conforme a norma ISO 9283 \cite{iso9283}, os testes de desempenho devem ser realizados em múltiplos pontos do espaço de trabalho, com número estatisticamente significativo de repetições e sob condições ambientais controladas. Portanto, a adoção desses critérios normalizados como base para a comparação experimental entre o protótipo EB-15 e o robô industrial UR-10 garante a objetividade e a reprodutibilidade das conclusões obtidas neste trabalho.
